\chapter{Τα Δομικά Στοιχεία της Υλοποίησης}
\label{ch:building-blocks}

Το παρόν κεφάλαιο παρουσιάζει τα λογισμικά εργαλεία που αξιοποιήθηκαν για την
υλοποίηση του \textlatin{Plegma Data Space}. Κάθε εργαλείο αντιστοιχεί σε έναν
συγκεκριμένο ρόλο εντός της αρχιτεκτονικής και επιλέχθηκε λόγω της συμβατότητάς
του με το πλαίσιο \textlatin{i4Trust} και της ανοιχτής φύσης του κώδικά του.

% ------------------------------------------------------------------
\section{\textlatin{FIWARE Orion-LD Context Broker}}
\label{sec:orion-ld}
% ------------------------------------------------------------------

Ο \textlatin{Orion-LD Context Broker}~\cite{orion-ld2023} είναι ο κεντρικός
κόμβος αποθήκευσης και διαχείρισης δεδομένων του \textlatin{Data Space}. Αποτελεί
εξέλιξη του αρχικού \textlatin{Orion Context Broker} της \textlatin{FIWARE Foundation}
με πλήρη υποστήριξη του προτύπου \textlatin{NGSI-LD}, επιτρέποντας την αποθήκευση
και ανάκτηση οντοτήτων με σημασιολογικό περιεχόμενο.

\subsection{Ρόλος στην Αρχιτεκτονική}
\label{subsec:orion-role}

Ο \textlatin{Orion-LD} επιτελεί τον ρόλο του \textbf{αποθετηρίου δεδομένων}
(\textlatin{data store}). Δέχεται αιτήσεις \textlatin{NGSI-LD} για δημιουργία,
ανάγνωση, ενημέρωση και διαγραφή οντοτήτων, και αποθηκεύει τα δεδομένα σε μια
βάση δεδομένων \textlatin{MongoDB}. Στην παρούσα υλοποίηση, ο \textlatin{Orion-LD}
\textbf{δεν είναι απευθείας προσβάσιμος} από εξωτερικούς χρήστες: όλες οι αιτήσεις
δρομολογούνται μέσω του \textlatin{Kong API Gateway}, εξασφαλίζοντας ότι κάθε
πρόσβαση ελέγχεται από τους μηχανισμούς εξουσιοδότησης.

\subsection{Τεχνικά Χαρακτηριστικά}
\label{subsec:orion-tech}

\begin{itemize}
    \item \textbf{Εικόνα \textlatin{Docker}:} \texttt{quay.io/fiware/orion-ld}
    \item \textbf{Θύρα:} \texttt{1026} (εσωτερικά, εντός του \textlatin{Docker network})
    \item \textbf{Αποθήκευση:} \textlatin{MongoDB} (βάση δεδομένων εγγράφων)
    \item \textbf{Πρωτόκολλο:} \textlatin{NGSI-LD v1} (\textlatin{ETSI GS CIM 009})
\end{itemize}

% ------------------------------------------------------------------
\section{\textlatin{Keyrock Identity Manager}}
\label{sec:keyrock}
% ------------------------------------------------------------------

Ο \textlatin{Keyrock}~\cite{keyrock2023} είναι ένας διαχειριστής ταυτότητας
(\textlatin{Identity Manager}) της \textlatin{FIWARE Foundation} που στην παρούσα
υλοποίηση επιτελεί έναν διπλό ρόλο: αυτόν του \textbf{παρόχου ταυτότητας}
(\textlatin{Identity Provider}) και αυτόν του \textbf{Μητρώου Εξουσιοδότησης}
(\textlatin{Authorization Registry, AR}).

\subsection{Ρόλος ως \textlatin{Authorization Registry}}
\label{subsec:keyrock-ar}

Στο πλαίσιο \textlatin{iSHARE}, το \textlatin{Μητρώο Εξουσιοδότησης} είναι ο
φορέας που αποθηκεύει και εκδίδει τις πολιτικές εκχώρησης εξουσιοδότησης
(\textlatin{delegation policies}). Ο \textlatin{Keyrock} υλοποιεί τις ακόλουθες
διεπαφές \textlatin{iSHARE}:

\begin{itemize}
    \item \texttt{POST /oauth2/token}: Αποδέχεται ένα \textlatin{iSHARE JWT} και
    επιστρέφει ένα \textlatin{access token} που ενσωματώνει τα στοιχεία
    \textlatin{delegation evidence}.
    \item \texttt{POST /ar/policy}: Επιτρέπει τη δημιουργία νέων πολιτικών
    εξουσιοδότησης από τον πάροχο δεδομένων.
    \item \texttt{POST /ar/delegation}: Επιστρέφει τα στοιχεία εκχώρησης
    εξουσιοδότησης για έναν συγκεκριμένο συνδυασμό πόρου και υποκειμένου.
\end{itemize}

\subsection{Επικοινωνία με τον \textlatin{iSHARE Satellite}}
\label{subsec:keyrock-satellite}

Κατά την επαλήθευση ενός \textlatin{iSHARE JWT} από καταναλωτή, ο \textlatin{Keyrock}
επικοινωνεί αυτόματα με τον \textlatin{iSHARE Satellite} για να επαληθεύσει ότι
ο εκδότης του \textlatin{token} είναι εγγεγραμμένος και αξιόπιστος συμμετέχων
στο \textlatin{Data Space}, μέσω του \textlatin{endpoint} \texttt{GET /parties}.

\subsection{Τεχνικά Χαρακτηριστικά}
\label{subsec:keyrock-tech}

\begin{itemize}
    \item \textbf{Εικόνα \textlatin{Docker}:} \texttt{quay.io/fiware/idm}
    \item \textbf{Θύρα:} \texttt{3007}
    \item \textbf{Αποθήκευση:} \textlatin{MySQL}
    \item \textbf{Αναγνωριστικό \textlatin{EORI}:} \texttt{EU.EORI.NLPLEGMA}
    (λειτουργεί ως \textlatin{Authorization Registry} του παρόχου)
\end{itemize}

% ------------------------------------------------------------------
\section{\textlatin{iSHARE Satellite}}
\label{sec:satellite}
% ------------------------------------------------------------------

Ο \textlatin{iSHARE Satellite}~\cite{ishare-satellite2023} επιτελεί τον ρόλο
της \textbf{Αρχής Εμπιστοσύνης} (\textlatin{Trust Anchor}) εντός του
\textlatin{Data Space}. Διαχειρίζεται το μητρώο των εγκεκριμένων συμμετεχόντων
και επαληθεύει την εγκυρότητα των ψηφιακών τους πιστοποιητικών.

\subsection{Ρόλος στην Αρχιτεκτονική}
\label{subsec:satellite-role}

Κάθε φορά που ένας συμμετέχων παρουσιάζει ένα \textlatin{iSHARE JWT}, ο
παραλήπτης επικοινωνεί με τον \textlatin{Satellite} για να:
\begin{enumerate}
    \item Αποκτήσει ένα \textlatin{access token} για τον ίδιο, μέσω του
    \texttt{POST /token}.
    \item Λάβει τη λίστα εμπιστευμένων \textlatin{CA} (\textlatin{trusted\_list}),
    μέσω του \texttt{GET /trusted\_list}, ώστε να επαληθεύσει την αλυσίδα
    πιστοποιητικών του \textlatin{JWT}.
    \item Επαληθεύσει ότι ο αιτών εκπροσωπεί ενεργό συμμετέχοντα, μέσω του
    \texttt{GET /parties}.
\end{enumerate}

\subsection{Παραμετροποίηση}
\label{subsec:satellite-config}

Η παραμετροποίηση του \textlatin{Satellite} γίνεται μέσω του αρχείου
\texttt{satellite.yml}, το οποίο ορίζει:
\begin{itemize}
    \item Τo αναγνωριστικό και τα κρυπτογραφικά κλειδιά του \textlatin{Satellite}.
    \item Τη λίστα εμπιστευμένων \textlatin{CA} (\textlatin{trusted\_list}).
    \item Το μητρώο εγκεκριμένων συμμετεχόντων (\textlatin{parties}), με τα
    στοιχεία του παρόχου (\texttt{EU.EORI.NLPLEGMA}) και του καταναλωτή
    (\texttt{EU.EORI.NLPLEGMACONSUMER}).
\end{itemize}

\subsection{Τεχνικά Χαρακτηριστικά}
\label{subsec:satellite-tech}

\begin{itemize}
    \item \textbf{Εικόνα \textlatin{Docker}:} \texttt{fiware/ishare-satellite:2.0.2}
    \item \textbf{Θύρα:} \texttt{8085} (εξωτερικά), \texttt{8080} (εσωτερικά)
    \item \textbf{Υλοποίηση:} \textlatin{Python/Flask} (\textlatin{WSGI})
    \item \textbf{Αναγνωριστικό \textlatin{EORI}:} \texttt{EU.EORI.NLPLEGMA}
\end{itemize}

% ------------------------------------------------------------------
\section{\textlatin{Kong API Gateway}}
\label{sec:kong}
% ------------------------------------------------------------------

Το \textlatin{Kong API Gateway}~\cite{kong2023} επιτελεί τον ρόλο του
\textbf{Σημείου Επιβολής Πολιτικής} (\textlatin{Policy Enforcement Point, PEP}).
Είναι το μοναδικό σημείο εισόδου για όλες τις αιτήσεις δεδομένων και εξασφαλίζει
ότι καμία αίτηση δεν φτάνει στον \textlatin{Orion-LD} χωρίς να έχει ελεγχθεί
η εξουσιοδότησή της.

\subsection{Το \textlatin{Plugin ngsi-ishare-policies}}
\label{subsec:kong-plugin}

Το \textlatin{Kong} επεκτείνεται με το \textlatin{plugin}
\texttt{ngsi-ishare-policies}~\cite{kong-ngsi-plugin2023}, ειδικά σχεδιασμένο
για \textlatin{Data Spaces} βάσει \textlatin{iSHARE}. Το \textlatin{plugin} αυτό,
κατά την παραλαβή κάθε αιτήσεως:

\begin{enumerate}
    \item Εξάγει το \textlatin{Bearer token} από την κεφαλίδα
    \texttt{Authorization} της αίτησης.
    \item Επαληθεύει την υπογραφή του \textlatin{token} χρησιμοποιώντας την
    αλυσίδα πιστοποιητικών \textlatin{X.509} από το πεδίο \texttt{x5c} και
    τη λίστα εμπιστευμένων \textlatin{CA} από τον \textlatin{Satellite}.
    \item Ελέγχει αν ο εκδότης (\texttt{iss}) είναι ο ίδιος ο πάροχος
    (απευθείας πρόσβαση) ή τρίτος οργανισμός (απαιτείται \textlatin{delegation
    evidence}).
    \item Σε περίπτωση τρίτου, ελέγχει το \textlatin{delegationEvidence}
    εντός του \textlatin{token} για να εξακριβώσει αν η αιτούμενη ενέργεια
    επιτρέπεται.
    \item Προωθεί την αίτηση στον \textlatin{Orion-LD} μόνο εφόσον η
    εξουσιοδότηση είναι έγκυρη.
\end{enumerate}

\subsection{Παραμετροποίηση μέσω \texttt{kong.yml}}
\label{subsec:kong-config}

Η παραμετροποίηση του \textlatin{Kong} γίνεται μέσω δηλωτικού αρχείου
\texttt{kong.yml} (\textlatin{declarative configuration}), χωρίς τη χρήση
βάσης δεδομένων (\textlatin{DB-less mode}). Το αρχείο ορίζει:
\begin{itemize}
    \item Την υπηρεσία-στόχο (\textlatin{upstream service}): \texttt{http://fiware-orion-ld:1026}
    \item Τα \textlatin{routes} που θα δέχεται: \texttt{/ngsi-ld}
    \item Τις παραμέτρους του \textlatin{plugin}: τα κλειδιά και πιστοποιητικά
    του παρόχου για υπογραφή \textlatin{JWT}, τα \textlatin{endpoints} του
    \textlatin{Satellite} και του \textlatin{Keyrock AR}.
\end{itemize}

\subsection{Τεχνικά Χαρακτηριστικά}
\label{subsec:kong-tech}

\begin{itemize}
    \item \textbf{Εικόνα \textlatin{Docker}:} \texttt{quay.io/fiware/kong:0.5.4}
    \item \textbf{Θύρα:} \texttt{8000} (\textlatin{proxy}), \texttt{8001}
    (\textlatin{admin API})
    \item \textbf{Λειτουργία:} \textlatin{DB-less} (δηλωτική παραμετροποίηση)
    \item \textbf{Αναγνωριστικό \textlatin{EORI}:} \texttt{EU.EORI.NLPLEGMA}
\end{itemize}

% ------------------------------------------------------------------
\section{\textlatin{DSBA Policy Decision Point (DSBA-PDP)}}
\label{sec:dsba-pdp}
% ------------------------------------------------------------------

Ο \textlatin{DSBA-PDP}~\cite{dsba-pdp2023} είναι ο \textbf{κόμβος λήψης
αποφάσεων πρόσβασης} (\textlatin{Policy Decision Point, PDP}) της αρχιτεκτονικής.
Αξιολογεί τις αιτήσεις πρόσβασης σε σχέση με τις ισχύουσες πολιτικές
\textlatin{iSHARE} και επιστρέφει απόφαση \textlatin{Permit} ή \textlatin{Deny}.

\subsection{Ρόλος στην Αρχιτεκτονική}
\label{subsec:pdp-role}

Στην παρούσα υλοποίηση, ο \textlatin{DSBA-PDP} ενεργοποιείται έμμεσα μέσω του
\textlatin{plugin} \texttt{ngsi-ishare-policies} του \textlatin{Kong}. Η λογική
αξιολόγησης των πολιτικών \textlatin{iSHARE} (\textlatin{delegation evidence})
υλοποιείται εντός του \textlatin{plugin}, ενώ ο \textlatin{DSBA-PDP} παρέχει
πρόσθετες δυνατότητες αξιολόγησης πολιτικών για πιο σύνθετα σενάρια.

\subsection{Τεχνικά Χαρακτηριστικά}
\label{subsec:pdp-tech}

\begin{itemize}
    \item \textbf{Εικόνα \textlatin{Docker}:} \texttt{quay.io/fiware/dsba-pdp:latest}
    \item \textbf{Θύρα:} \texttt{8080}
    \item \textbf{Παραμετροποίηση:} Αναγνωριστικό \textlatin{EORI} παρόχου,
    διαδρομές πιστοποιητικών, \textlatin{endpoint} \textlatin{Satellite} και
    \textlatin{Keyrock AR}
\end{itemize}

% ------------------------------------------------------------------
\section{Βάσεις Δεδομένων}
\label{sec:databases}
% ------------------------------------------------------------------

Η υλοποίηση χρησιμοποιεί δύο βάσεις δεδομένων, καθεμία εξυπηρετώντας
διαφορετικές ανάγκες:

\subsection{\textlatin{MongoDB}}
\label{subsec:mongodb}

Η \textlatin{MongoDB} χρησιμοποιείται από τον \textlatin{Orion-LD} για την
αποθήκευση των οντοτήτων \textlatin{NGSI-LD}. Ως βάση δεδομένων εγγράφων
(\textlatin{document store}), η \textlatin{MongoDB} είναι ιδανική για την
αποθήκευση οντοτήτων \textlatin{JSON-LD}, καθώς δεν απαιτεί προκαθορισμένο
σχήμα (\textlatin{schema-less}).

\subsection{\textlatin{MySQL}}
\label{subsec:mysql}

Η \textlatin{MySQL} χρησιμοποιείται από τον \textlatin{Keyrock} για την
αποθήκευση χρηστών, εφαρμογών και πολιτικών εξουσιοδότησης. Κατά την εκκίνηση
του συστήματος, ο \textlatin{Keyrock} δημιουργεί αυτόματα τις απαραίτητες
εγγραφές για τους δύο συμμετέχοντες (\texttt{EU.EORI.NLPLEGMA} και
\texttt{EU.EORI.NLPLEGMACONSUMER}) μέσω του μηχανισμού
\textlatin{extparticipant}.

% ------------------------------------------------------------------
\section{Συγκριτική Αξιολόγηση Εργαλείων}
\label{sec:comparison}
% ------------------------------------------------------------------

Αξίζει να σημειωθεί ότι η \textlatin{FIWARE Foundation} παρέχει επίσης έναν
\textlatin{Data Space Connector}, ο οποίος αποτελεί μια εναλλακτική, πιο
ολοκληρωμένη προσέγγιση υλοποίησης \textlatin{Data Spaces} βάσει νεότερης
αρχιτεκτονικής (\textlatin{OID4VC}/\textlatin{Verifiable Credentials}). Η
παρούσα εργασία ακολουθεί το παλαιότερο, αλλά πλήρως λειτουργικό, μοντέλο
\textlatin{i4Trust}/\textlatin{iSHARE}, το οποίο αντιστοιχεί ρητά στις
απαιτήσεις της εκφώνησης και παρέχει μια σαφή, ευέλικτη αρχιτεκτονική με
ανεξάρτητα, ελεγχόμενα συστατικά.

Τα πλεονεκτήματα της επιλεγμένης προσέγγισης είναι:
\begin{itemize}
    \item \textbf{Ανοιχτός κώδικας:} Όλα τα εργαλεία διατίθενται υπό ανοιχτές
    άδειες χρήσης.
    \item \textbf{Ευελιξία:} Κάθε συστατικό μπορεί να αντικατασταθεί ανεξάρτητα.
    \item \textbf{Διαφάνεια:} Η λογική εξουσιοδότησης είναι ορατή και ελέγξιμη
    σε κάθε βήμα της ροής.
    \item \textbf{Δυνατότητα τοπικής ανάπτυξης:} Το σύστημα μπορεί να
    εκτελεστεί πλήρως σε τοπικό περιβάλλον με \textlatin{Docker Compose},
    χωρίς εξάρτηση από εξωτερικές υπηρεσίες.
\end{itemize}

% ------------------------------------------------------------------
\section{Σύνοψη}
\label{sec:ch4-summary}
% ------------------------------------------------------------------

Στο παρόν κεφάλαιο παρουσιάστηκαν τα πέντε κύρια συστατικά της υλοποίησης:

\begin{itemize}
    \item \textbf{\textlatin{Orion-LD}}: Αποθετήριο δεδομένων \textlatin{NGSI-LD}.
    \item \textbf{\textlatin{Keyrock}}: Διαχειριστής ταυτότητας και Μητρώο
    Εξουσιοδότησης.
    \item \textbf{\textlatin{iSHARE Satellite}}: Αρχή εμπιστοσύνης και μητρώο
    συμμετεχόντων.
    \item \textbf{\textlatin{Kong API Gateway}}: Σημείο επιβολής πολιτικής.
    \item \textbf{\textlatin{DSBA-PDP}}: Κόμβος λήψης αποφάσεων πρόσβασης.
\end{itemize}

Στο επόμενο κεφάλαιο παρουσιάζεται η αρχιτεκτονική που προκύπτει από τη
σύνθεση αυτών των συστατικών, με αναλυτικά διαγράμματα και περιγραφή των
δικτυακών διασυνδέσεων τους.
